% GENERATED FILE. Edit canonical lessons, then run npm run generate:books.
% canonical-source-hash: fnv1a64:51e5db5ac3040066
% canonical-lessons: ML-C123-notpossible, ML-C123-notlikethat, ML-C123-maybe, ML-C123-other, ML-C123-wrong

\chapter{Maybe, and Not}
\label{ch:ml-maybe-and-not}
\hlchaptermodality{123}

\begin{chapteropening}
\textbf{By the end of this chapter:} \emph{I can say it is not possible, not like that, maybe, other, and wrong.}

Complete the last lesson of chapter 123: i can say it is not possible, not like that, maybe, other, and wrong.
\end{chapteropening}

\section[paṟṟilla]{\ml{പറ്റില്ല} (paṟṟilla) — it's not possible, I can't}

\label{lesson:ML-C123-notpossible}



\begin{quote}
Before the new one: say the Malayalam for the same, then the Malayalam for oh, I see.
\end{quote}

\subsection*{You'll want to know: \ml{പറ്റില്ല}}
\textbf{\ml{പറ്റില്ല}} (\emph{paṟṟilla}) — \textquotedblleft{}it's not possible, (I) can't\textquotedblright{}. A firm, everyday no.

\begin{grammarlens}[title={What you've built}]
The polite but firm no.
\end{grammarlens}

\subsection*{Your turn}
\begin{itemize}
\raggedright
  \item \emph{Say it:} \emph{paṟṟilla}
  \item \emph{Say it:} \emph{paṟṟilla}, once more
  \item \emph{Say it:} say \emph{paṟṟilla}
  \item \emph{Recall:} say the Malayalam for the same, then the Malayalam for oh, I see, then say \emph{paṟṟilla} again
\end{itemize}

\begin{tcolorbox}[breakable,colback=teal!4,colframe=teal!35!black,title={Before you move on}]
What does \ml{പറ്റില്ല} mean? (\textquotedblleft{}It's not possible, I can't\textquotedblright{}.) Say it once more.
\end{tcolorbox}
\section[aṅṅaneyalla]{\ml{അങ്ങനെയല്ല} (aṅṅaneyalla) — it's not like that}

\label{lesson:ML-C123-notlikethat}



\begin{quote}
Before the new one: say the Malayalam for oh, I see, then the Malayalam for not possible.
\end{quote}

\subsection*{You'll want to know: \ml{അങ്ങനെയല്ല}}
\textbf{\ml{അങ്ങനെയല്ല}} (\emph{aṅṅaneyalla}) — \textquotedblleft{}it's not like that\textquotedblright{}: \textbf{\ml{അങ്ങനെ}} (\emph{aṅṅane}), \textquotedblleft{}like that\textquotedblright{}, with \textbf{\ml{അല്ല}} (\emph{alla}), \textquotedblleft{}is not\textquotedblright{}.

\begin{grammarlens}[title={What you've built}]
Two known pieces, one correction.
\end{grammarlens}

\subsection*{Your turn}
\begin{itemize}
\raggedright
  \item \emph{Say it:} \emph{aṅṅaneyalla}
  \item \emph{Say it:} \emph{aṅṅaneyalla}, once more
  \item \emph{Say it:} say \emph{aṅṅaneyalla}
  \item \emph{Recall:} say the Malayalam for oh, I see, then the Malayalam for not possible, then say \emph{aṅṅaneyalla} again
\end{itemize}

\begin{tcolorbox}[breakable,colback=teal!4,colframe=teal!35!black,title={Before you move on}]
What does \ml{അങ്ങനെയല്ല} mean? (\textquotedblleft{}It's not like that\textquotedblright{}.) Say it once more.
\end{tcolorbox}
\section[āyirikkāṁ]{\ml{ആയിരിക്കാം} (āyirikkāṁ) — maybe, it may be}

\label{lesson:ML-C123-maybe}



\begin{quote}
Before the new one: say the Malayalam for not possible, then the Malayalam for not like that.
\end{quote}

\subsection*{You'll want to know: \ml{ആയിരിക്കാം}}
\textbf{\ml{ആയിരിക്കാം}} (\emph{āyirikkāṁ}) — \textquotedblleft{}maybe, it could be\textquotedblright{}. \textbf{\ml{ഒരുപക്ഷേ}} (\emph{orupakṣē}), \textquotedblleft{}perhaps\textquotedblright{}, often comes before it.

\begin{grammarlens}[title={What you've built}]
A guess, softly put.
\end{grammarlens}

\subsection*{Your turn}
\begin{itemize}
\raggedright
  \item \emph{Say it:} \emph{āyirikkāṁ}
  \item \emph{Say it:} \emph{āyirikkāṁ}, once more
  \item \emph{Say it:} say \emph{orupakṣē āyirikkāṁ}
  \item \emph{Recall:} say the Malayalam for not possible, then the Malayalam for not like that, then say \emph{āyirikkāṁ} again
\end{itemize}

\begin{tcolorbox}[breakable,colback=teal!4,colframe=teal!35!black,title={Before you move on}]
What does \ml{ആയിരിക്കാം} mean? (\textquotedblleft{}Maybe, it may be\textquotedblright{}.) Say it once more.
\end{tcolorbox}
\section[vēṟe]{\ml{വേറെ} (vēṟe) — other, different}

\label{lesson:ML-C123-other}



\begin{quote}
Before the new one: say the Malayalam for not like that, then the Malayalam for maybe.
\end{quote}

\subsection*{You'll want to know: \ml{വേറെ}}
\textbf{\ml{വേറെ}} (\emph{vēṟe}) — \textquotedblleft{}other, different\textquotedblright{}. \textbf{\ml{വേറെ} \ml{വേണം}} (\emph{vēṟe vēṇaṁ}) — \textquotedblleft{}I want a different one\textquotedblright{}.

\begin{grammarlens}[title={What you've built}]
When the first choice is wrong.
\end{grammarlens}

\subsection*{Your turn}
\begin{itemize}
\raggedright
  \item \emph{Say it:} \emph{vēṟe}
  \item \emph{Say it:} \emph{vēṟe}, once more
  \item \emph{Say it:} say \emph{vēṟe vēṇaṁ}
  \item \emph{Recall:} say the Malayalam for not like that, then the Malayalam for maybe, then say \emph{vēṟe} again
\end{itemize}

\begin{tcolorbox}[breakable,colback=teal!4,colframe=teal!35!black,title={Before you move on}]
What does \ml{വേറെ} mean? (\textquotedblleft{}Other, different\textquotedblright{}.) Say it once more.
\end{tcolorbox}
\section[teṟṟŭ]{\ml{തെറ്റ്} (teṟṟŭ) — wrong; a mistake}

\label{lesson:ML-C123-wrong}



\begin{quote}
Before the new one: say the Malayalam for maybe, then the Malayalam for other.
\end{quote}

\subsection*{You'll want to know: \ml{തെറ്റ്}}
\textbf{\ml{തെറ്റ്}} (\emph{teṟṟŭ}) — \textquotedblleft{}wrong\textquotedblright{}, and also \textquotedblleft{}a mistake\textquotedblright{}. It is the honest partner of \textbf{\ml{ശരി}}, \textquotedblleft{}right\textquotedblright{}.

\begin{grammarlens}[title={What you've built}]
The opposite of śari.
\end{grammarlens}

\subsection*{Your turn}
\begin{itemize}
\raggedright
  \item \emph{Say it:} \emph{teṟṟŭ}
  \item \emph{Say it:} \emph{teṟṟŭ}, once more
  \item \emph{Say it:} say \emph{teṟṟŭ}
  \item \emph{Recall:} say the Malayalam for maybe, then the Malayalam for other, then say \emph{teṟṟŭ} again
\end{itemize}

\begin{tcolorbox}[breakable,colback=teal!4,colframe=teal!35!black,title={Before you move on}]
What does \ml{തെറ്റ്} mean? (\textquotedblleft{}Wrong; a mistake\textquotedblright{}.) Say it once more.
\end{tcolorbox}
